%==============================================================================
%  CHAPTER 4 — SPRINT 1: FOUNDATION (AUTHENTICATION & PROJECT MANAGEMENT)
%==============================================================================
\chapter{Sprint 1 --- Foundation: Authentication and Project Management}
\label{chap:sprint1}

\section*{Introduction}
The first sprint lays the foundation of the platform: a secure authentication
system and the project entity around which all later work is organized. This
chapter presents the sprint backlog, the analysis and conception, the realization,
and the sprint review.

% ---------------------------------------------------------------------------
\section{Sprint Backlog}
\begin{table}[H]\centering
\caption{Sprint~1 backlog}\label{tab:backlog-s1}
\begin{tabularx}{\textwidth}{p{0.6cm} X p{2cm}}
\toprule
\textbf{\#} & \textbf{Task} & \textbf{Pts}\\
\midrule
1 & Design the user and project data model & 2\\
1 & Implement signup / login / refresh with JWT & 3\\
1 & Build the login and signup user interface & 2\\
2 & Implement project CRUD with ownership checks & 3\\
2 & Bind each project to its own tracking experiment & 1\\
2 & Build the projects list and detail interface & 2\\
\bottomrule
\end{tabularx}
\end{table}

% ---------------------------------------------------------------------------
\section{Analysis and Conception}

\subsection{Use Cases}
This sprint covers \emph{Sign up}, \emph{Sign in}, \emph{Refresh session}, and
\emph{Manage projects} (create, list, update, delete).

\subsection{Sequence Diagram: Authentication}
% >>> USER: render diagrams/seq_login.puml -> images/seq_login.pdf
\begin{figure}[H]\centering
  % \includegraphics[width=0.85\textwidth]{seq_login.pdf}
  \fbox{\parbox[c][5cm][c]{0.85\textwidth}{\centering\itshape
    Placeholder --- insert \texttt{images/seq\_login.pdf}\\
    (render \texttt{diagrams/seq\_login.puml})}}
  \caption{Sequence diagram of the login flow}
  \label{fig:seq-login}
\end{figure}

\subsection{Data Model Increment}
The sprint introduces two entities: \texttt{User} (unique email, hashed password,
full name) and \texttt{Project} (name, owner reference, tracking-experiment
identifier).
% >>> USER: render diagrams/class_sprint1.puml -> images/class_sprint1.pdf
\begin{figure}[H]\centering
  % \includegraphics[width=0.6\textwidth]{class_sprint1.pdf}
  \fbox{\parbox[c][4cm][c]{0.6\textwidth}{\centering\itshape
    Placeholder --- \texttt{images/class\_sprint1.pdf}}}
  \caption{Class diagram increment after Sprint~1}
  \label{fig:class-s1}
\end{figure}

% ---------------------------------------------------------------------------
\section{Realization}

\subsection{Authentication with JWT}
Authentication is stateless and based on \gls{jwt}. On login, the backend verifies
the password against its stored hash and issues a short-lived \emph{access} token
and a longer-lived \emph{refresh} token. Protected endpoints depend on a function
that decodes the access token and resolves the current user; the refresh endpoint
exchanges a valid refresh token for a new access token without re-entering
credentials. Passwords are never stored in clear text, only a salted hash is
persisted.

\begin{lstlisting}[language=Python,caption={Resolving the current user from a JWT (backend dependency).},label={lst:get-current-user}]
async def get_current_user(token: str = Depends(oauth2_scheme),
                           db: AsyncSession = Depends(get_db)) -> User:
    payload = decode_access_token(token)        # raises 401 if invalid/expired
    user = await db.get(User, payload["sub"])
    if user is None:
        raise HTTPException(status_code=401, detail="Invalid credentials")
    return user
\end{lstlisting}

\subsection{Project Management}
A project is the organizing unit of the platform. Creating one also provisions a
dedicated experiment in the tracking server, so that every run under the project
is grouped automatically. All project endpoints enforce an ownership check: a user
can only see and modify their own projects.

\subsection{User Interface}
% >>> USER: add screenshots to images/ (login + projects list) and uncomment.
% \begin{figure}[H]\centering
%   \includegraphics[width=\textwidth]{screenshot_login.png}
%   \caption{Login screen}\label{fig:ui-login}
% \end{figure}
The frontend provides a login/signup screen and a projects view from which the
user creates and opens projects.

% ---------------------------------------------------------------------------
\section{Sprint Review and Retrospective}
At the end of Sprint~1 the platform supports secure authentication and full
project management, the foundation every later sprint builds upon. The burndown
chart below tracks remaining effort across the sprint.
% >>> USER: replace the y-values with your real remaining points per day.
\begin{figure}[H]\centering
\begin{tikzpicture}
\begin{axis}[
    width=0.8\textwidth, height=6cm,
    xlabel={Day}, ylabel={Remaining points},
    xmin=0, xmax=10, ymin=0, ymax=14,
    legend pos=north east, grid=both,
]
\addplot[dashed,gray] coordinates {(0,13)(9,0)};         % ideal
\addplot[thick,primary,mark=*] coordinates {              % actual
  (0,13)(1,13)(2,11)(3,9)(4,8)(5,6)(6,5)(7,3)(8,1)(9,0)};
\legend{Ideal, Actual}
\end{axis}
\end{tikzpicture}
\caption{Sprint~1 burndown chart}\label{fig:burndown-s1}
\end{figure}

\noindent\textbf{Retrospective.} % >>> USER: 2--3 lines of what went well / to improve.
What went well: the authentication and project foundations were delivered on
schedule. To improve: \todo{your note}.
